\documentclass[a4paper,11pt]{scrartcl}

\usepackage[T1]{fontenc}
\usepackage[utf8]{inputenc}
\usepackage[ngerman]{babel}
\usepackage{lmodern}
\usepackage{amsmath}
\usepackage{amssymb}
\usepackage[a4paper,margin=2.5cm]{geometry}
\usepackage{enumitem}
\usepackage{booktabs}
\usepackage[hidelinks]{hyperref}
\usepackage{catppuccin-dark}

\newcounter{aufgabe}
\newcommand{\Aufgabe}[2]{%
  \refstepcounter{aufgabe}%
  \subsection*{Aufgabe~\theaufgabe: #1 \hfill (#2~Punkte)}%
  \label{auf:\theaufgabe}}

\title{Übungsblatt 3: Error Analysis (Fehleranalyse)}
\subtitle{Begleitmaterial zur Vorlesung \glqq Numerical Methods\grqq{} (Prof.\ Dr.-Ing.\ Mark Schutera)}
\author{Konditionierung, Stabilität, Konsistenz und Konvergenz\\(Skript S.~25--34)}
\date{\today}

\begin{document}

\maketitle

\noindent\textbf{Gesamt: 50 Punkte.}\quad
Hilfsmittel: Taschenrechner bzw.\ die angegebenen Wertetabellen. Runden Sie, sofern nicht anders angegeben, auf fünf Nachkommastellen (Endergebnisse vier Nachkommastellen genügen). Die Aufgaben steigen im Schwierigkeitsgrad an: Aufgaben 1--3 prüfen das Begriffsverständnis, Aufgaben 4--6 sind Rechenaufgaben analog zum Skript (mit eigenem Beispiel), Aufgaben 7--8 sind Transferaufgaben. Die Gleichungsnummern (13)--(19) beziehen sich auf das Skript.

\bigskip
\hrule

\section*{Aufgaben}

\Aufgabe{Hut- und Tilde-Notation (Hat versus Tilde)}{4}

Im Skript bezeichnet $f(\cdot)$ die exakte (analytische) Lösung eines Problems und $\hat{f}(\cdot)$ die numerische Methode (den Algorithmus), die eine Approximation liefert. $x$ steht für die exakten Eingabedaten, $\tilde{x}$ für die tatsächlich verwendeten, möglicherweise gestörten Eingabedaten.

\begin{enumerate}[label=(\alph*)]
  \item Erklären Sie in jeweils einem Satz, was die vier Symbole $f$, $\hat{f}$, $x$ und $\tilde{x}$ bedeuten und wodurch eine Störung von $x$ zu $\tilde{x}$ in der Praxis typischerweise entsteht. \hfill (2~P)
  \item Übersetzen Sie die vier Ausdrücke
  \[ f(x), \qquad f(\tilde{x}), \qquad \hat{f}(x), \qquad \hat{f}(\tilde{x}) \]
  jeweils in Worte. Welcher der vier Ausdrücke liest sich als \glqq numerische Methode, ausgewertet auf gestörten Eingabedaten\grqq{} (numerical method evaluated on perturbed input data)? \hfill (2~P)
\end{enumerate}

\schreibraum[7]
\Aufgabe{Die vier Grundbegriffe unterscheiden und zuordnen}{6}

Gegeben sind vier verbale Beschreibungen:

\begin{enumerate}[label=\textbf{\Alph*.}]
  \item Wie sensitiv reagiert die Lösung eines \emph{Problems} auf kleine Änderungen der Eingabedaten?
  \item Kleine Fehler in der Eingabe oder in Zwischenschritten führen nicht zu unverhältnismäßig großen Fehlern in der Ausgabe der \emph{Methode}; der Ausgabefehler bleibt durch ein konstantes Vielfaches des Eingabefehlers beschränkt.
  \item Wie gut approximiert die numerische Lösung die exakte Lösung des Originalproblems, wenn beide auf denselben, \emph{ungestörten} Eingabedaten ausgewertet werden?
  \item Die Approximation erreicht einen spezifischen, stabilen Grenzwert: Sowohl die Annäherung an die echte Lösung als auch die Abweichungen in den Daten werden durch eine Fehlerschranke kontrolliert.
\end{enumerate}

und vier Formeln:
\begin{align*}
\text{(i)}\quad & \left|\hat{f}(x) - f(x)\right| \leq c
& \text{(ii)}\quad & \left|\hat{f}(\tilde{x}) - f(x)\right| \to \lim \\
\text{(iii)}\quad & \kappa = \left|f(x) - f(\tilde{x})\right|
& \text{(iv)}\quad & \left|\hat{f}(\tilde{x}) - \hat{f}(x)\right| \leq s \cdot |\tilde{x} - x|
\end{align*}

\begin{enumerate}[label=(\alph*)]
  \item Ordnen Sie jeder Beschreibung A--D den passenden Fachbegriff --- Konditionierung (conditioning), Stabilität (stability), Konsistenz (consistency), Konvergenz (convergence) --- sowie die passende Formel (i)--(iv) zu. \hfill (4~P)
  \item Welche der vier Begriffe sind Eigenschaften des \emph{Problems}, welche Eigenschaften der \emph{Methode}? Begründen Sie kurz anhand der in den Formeln auftretenden Symbole. \hfill (1~P)
  \item Begründen Sie ebenfalls anhand der auftretenden Symbole, warum man sagt: \emph{Konvergenz kombiniert die Ideen von Stabilität und Konsistenz.} \hfill (1~P)
\end{enumerate}

\schreibraum[9]
\Aufgabe{Nicht verschwindende Konvergenz und systematischer Fehler (Bias)}{4}

Eine numerische Methode $\hat{f}_h$ wird auf gestörten Eingabedaten $\tilde{x}$ ausgewertet. Für immer kleinere Schrittweiten $h$ misst man die Gesamtfehler $\left|\hat{f}_h(\tilde{x}) - f(x)\right|$:

\begin{center}
\begin{tabular}{lcccccc}
\toprule
$h$ & $0.4$ & $0.2$ & $0.1$ & $0.05$ & $0.025$ & $0.0125$ \\
\midrule
$\left|\hat{f}_h(\tilde{x}) - f(x)\right|$ & $0.080$ & $0.045$ & $0.034$ & $0.0315$ & $0.0311$ & $0.0310$ \\
\bottomrule
\end{tabular}
\end{center}

\begin{enumerate}[label=(\alph*)]
  \item Konvergiert die Methode im Sinne von Gl.~(16)? Gegen welche Fehlerschranke (ungefähr)? \hfill (2~P)
  \item Konvergenz zielt üblicherweise auf die Fehlerschranke null, also die exakte Lösung $f(x)$. Was zeigt es nach dem Skript an, wenn die Methode --- wie hier --- gegen einen von $f(x)$ verschiedenen Wert konvergiert? \hfill (1~P)
  \item Drücken Sie die Beobachtung in den Begriffen des maschinellen Lernens \emph{Bias} und \emph{Varianz} aus: Welcher Anteil des Fehlerverhaltens entspricht dem Bias, welcher der Varianz? \hfill (1~P)
\end{enumerate}

\subsection*{Gemeinsames Szenario für die Aufgaben 4--6}

Gesucht ist das Minimum der Funktion $f(x) = \cos(x)$ auf dem Intervall $[0, 4]$. Die exakte Lösung ist $f(x^*) = -1$ bei $x^* = \pi \approx 3.14159$. Die numerische Methode diskretisiert das Intervall mit Schrittweite $h$ (Gitterpunkte $x_i = i \cdot h$) und wertet stückweise konstant am \emph{nächstgelegenen} Gitterpunkt aus:
\[ \hat{f}_h(x) = f(x_i), \qquad x_i \text{ ist der zu } x \text{ nächstgelegene Gitterpunkt.} \]
Betrachtet werden die Schrittweiten $h = 1$ und $h = 0.2$, also die Gitter $\{0, 1, 2, 3, 4\}$ bzw.\ $\{0,\, 0.2,\, 0.4,\, \ldots,\, 4\}$. Die Eingabe $x$ ist durch Messfehler gestört: $\tilde{x} = x \pm 0.25$. Aus Symmetriegründen betrachten wir --- wie im Skript --- jeweils nur den Fall $\tilde{x} = x + 0.25$, sofern nichts anderes verlangt ist.

\medskip
\noindent\textbf{Wertetabelle} (auf fünf Nachkommastellen gerundet):
\begin{center}
\begin{tabular}{lcccccc}
\toprule
$x$ & $1.25$ & $1.5$ & $1.75$ & $2.0$ & $2.2$ & $2.25$ \\
$\cos(x)$ & $0.31532$ & $0.07074$ & $-0.17825$ & $-0.41615$ & $-0.58850$ & $-0.62817$ \\
\midrule
$x$ & $2.75$ & $3.0$ & $3.2$ & $3.25$ & $3.4$ & \\
$\cos(x)$ & $-0.92430$ & $-0.98999$ & $-0.99829$ & $-0.99413$ & $-0.96680$ & \\
\bottomrule
\end{tabular}
\end{center}

\schreibraum[7]
\Aufgabe{Konditionierung: Konditionszahl $\kappa$ nach Gl.~(13)}{6}

Die Konditionierung beschreibt die Sensitivität des \emph{Problems} selbst und ist unabhängig von der numerischen Methode. Quantifizieren Sie sie über die Konditionszahl (condition number) nach Gl.~(13),
\[ \kappa = \left| f(x) - f(\tilde{x}) \right|, \]
mit $f(x) = \cos(x)$ und der Störung $\tilde{x} = x \pm 0.25$ (hier \emph{beide} Vorzeichen).

\begin{enumerate}[label=(\alph*)]
  \item Berechnen Sie $\kappa(3, +0.25)$ und $\kappa(3, -0.25)$ am Punkt $x = 3$ (nahe der Minimalstelle $\pi$). \hfill (2~P)
  \item Berechnen Sie $\kappa(1.5, +0.25)$ und $\kappa(1.5, -0.25)$ am Punkt $x = 1.5$ (nahe der Nullstelle $\pi/2$). \hfill (2~P)
  \item Klassifizieren Sie beide Punkte als gut konditioniert (well-conditioned) bzw.\ schlecht konditioniert (ill-conditioned) und begründen Sie das Ergebnis über das Steigungsverhalten von $\cos(x)$ an den beiden Punkten. \hfill (2~P)
\end{enumerate}

\schreibraum[9]
\Aufgabe{Stabilität: Stabilitätskonstante $s$ nach Gl.~(14)/(17)}{7}

Die Stabilität hängt von der numerischen \emph{Methode} ab. Quantifizieren Sie die Konstante $s$ in der Stabilitätsungleichung (Gl.~(14), im Skript als Gl.~(17) wiederholt),
\[ \left| \hat{f}(\tilde{x}) - \hat{f}(x) \right| \leq s \cdot |\tilde{x} - x|, \]
am Punkt $x = 2$ mit $\tilde{x} = x + 0.25 = 2.25$.

\begin{enumerate}[label=(\alph*)]
  \item Für $h = 1$: Bestimmen Sie die zu $x$ und $\tilde{x}$ nächstgelegenen Gitterpunkte, werten Sie $\hat{f}_1$ aus und berechnen Sie $s_{h=1,\,x=2}$. \hfill (3~P)
  \item Für $h = 0.2$: ebenso für $s_{h=0.2,\,x=2}$. \hfill (3~P)
  \item Welche der beiden Diskretisierungen ist an diesem Punkt stabiler, und warum --- was \glqq sieht\grqq{} das feine Gitter, das das grobe Gitter nicht sieht? \hfill (1~P)
\end{enumerate}

\schreibraum[10]
\Aufgabe{Konsistenz nach Gl.~(15)/(18) und Konvergenz nach Gl.~(16)/(19)}{9}

\begin{enumerate}[label=(\alph*)]
  \item \textbf{Konsistenz:} Quantifizieren Sie die Konsistenzkonstante $c$ aus
  \[ \left| \hat{f}(x) - f(x) \right| \leq c \]
  an der exakten Minimalstelle $x^* = \pi \approx 3.14159$ (mit $f(x^*) = -1$) für beide Schrittweiten, d.\,h.\ bestimmen Sie untere Schranken $c_1$ und $c_{0.2}$ analog zur Rechnung im Skript. \hfill (3~P)
  \item \textbf{Konvergenz:} Quantifizieren Sie den Gesamtfehler
  \[ \left| \hat{f}(\tilde{x}) - f(x^*) \right| \]
  für die gestörte Eingabe $\tilde{x} = x^* + 0.25 \approx 3.39159$, ebenfalls für beide Schrittweiten. \hfill (4~P)
  \item Vergleichen Sie: Wie verändern sich Konsistenz und Konvergenz beim Übergang von $h = 1$ zu $h = 0.2$? Formulieren Sie den (auf den ersten Blick überraschenden) Befund in einem Satz. \hfill (2~P)
\end{enumerate}

\schreibraum[12]
\Aufgabe{Transfer --- Beweisaufgabe: Stabilität wird zur Konditionierung für $h \to 0$}{8}

Das Skript stellt in einer Randnotiz (S.~31) die Behauptung auf: \glqq Stability equals the conditioning of the system, for $h \to 0$. Show it.\grqq{} --- Die Stabilität entspricht im Grenzfall $h \to 0$ der Konditionierung des Systems. Zeigen Sie dies für die stückweise konstante Auswertung am nächstgelegenen Gitterpunkt und eine stetig differenzierbare Funktion $f$ mit beschränkter Ableitung, $L := \max_{x \in [a,b]} |f'(x)|$.

\begin{enumerate}[label=(\alph*)]
  \item Zeigen Sie zunächst: $\left| \hat{f}_h(x) - f(x) \right| \leq L \cdot \dfrac{h}{2}$ für alle $x \in [a,b]$, d.\,h.\ die Methode wird für $h \to 0$ punktweise exakt. \hfill (3~P)
  \item Folgern Sie daraus mit der Dreiecksungleichung:
  \[ \left| \hat{f}_h(\tilde{x}) - \hat{f}_h(x) \right| \;\xrightarrow{\;h \to 0\;}\; \left| f(\tilde{x}) - f(x) \right| = \kappa, \]
  und interpretieren Sie das Ergebnis: Was passiert mit der linken Seite der Stabilitätsungleichung (Gl.~14), und was bedeutet das für die kleinstmögliche Stabilitätskonstante $s$? \hfill (3~P)
  \item Überprüfen Sie die Aussage numerisch am Szenario der Aufgaben 4--6 (Punkt $x = 2$, $\tilde{x} = 2.25$): Vergleichen Sie die Werte $\left|\hat{f}_h(\tilde{x}) - \hat{f}_h(x)\right|$ für $h = 1$ und $h = 0.2$ (aus Aufgabe~5) sowie für $h = 0.05$ mit der Konditionszahl $\kappa(2, +0.25)$. \hfill (2~P)
\end{enumerate}

\schreibraum[11]
\Aufgabe{Transfer --- Zusammenspiel der Konzepte und Grenzen der $h$-Verfeinerung}{6}

\begin{enumerate}[label=(\alph*)]
  \item Fassen Sie anhand Ihrer Ergebnisse aus den Aufgaben 5 und 6 zusammen, wie Konsistenz, Stabilität und Konvergenz beim Übergang $h = 1 \to h = 0.2$ zusammenspielen. Warum garantiert eine bessere Konsistenz allein noch keine bessere Konvergenz? \hfill (2~P)
  \item Bestimmen Sie mit Aufgabe~7 den Grenzwert des Konvergenzfehlers $\left|\hat{f}_h(\tilde{x}) - f(x^*)\right|$ für $h \to 0$ im Szenario aus Aufgabe~6\,(b). \emph{Hinweis:} $\cos(\pi + 0.25) = -0.96891$. Wie nennt man dieses Verhalten (vgl.\ Aufgabe~3)? \hfill (2~P)
  \item Nennen und erläutern Sie zwei voneinander unabhängige Gründe, warum sich die Genauigkeit der numerischen Lösung \emph{nicht} beliebig durch weiteres Verkleinern von $h$ steigern lässt. Gehen Sie bei einem der beiden Gründe explizit auf die Gleitkommadarstellung (floating-point representation) ein. \hfill (2~P)
\end{enumerate}

\bigskip
\hrule
\medskip
\noindent\emph{Summe: 50 Punkte. Viel Erfolg!}

\schreibraum[9]
\newpage

\section*{Lösungen --- erst nach Bearbeitung lesen!}

\subsection*{Lösung zu Aufgabe~\ref{auf:1}}

\begin{enumerate}[label=(\alph*)]
  \item
  \begin{itemize}[nosep]
    \item $f$: die exakte (analytische) Lösung bzw.\ das exakte Modell des Problems.
    \item $\hat{f}$: die numerische Methode (der Algorithmus), die eine Approximation der exakten Lösung liefert --- der \emph{Hut} markiert die Methode.
    \item $x$: die exakten Eingabedaten.
    \item $\tilde{x}$: die tatsächlich verwendeten Eingabedaten --- die \emph{Tilde} markiert die gestörte Eingabe. Störungen entstehen in der Praxis typischerweise durch Messfehler (measurement error) oder Rundung (rounding) bei der Speicherung im Rechner.
  \end{itemize}
  \item
  \begin{itemize}[nosep]
    \item $f(x)$: exakte Lösung auf exakten Eingabedaten (die \glqq Wahrheit\grqq).
    \item $f(\tilde{x})$: exakte Lösung auf gestörten Eingabedaten (das Problem reagiert auf die Störung --- Gegenstand der Konditionierung).
    \item $\hat{f}(x)$: numerische Methode auf exakten Eingabedaten (reiner Methodenfehler --- Gegenstand der Konsistenz).
    \item $\hat{f}(\tilde{x})$: \textbf{numerische Methode, ausgewertet auf gestörten Eingabedaten} --- das ist der gesuchte Ausdruck (Merkhilfe des Skripts, Randnotiz S.~29: \glqq So $\hat{f}(\tilde{x})$ reads as: numerical method evaluated on perturbed input data.\grqq).
  \end{itemize}
\end{enumerate}

\subsection*{Lösung zu Aufgabe~\ref{auf:2}}

\begin{enumerate}[label=(\alph*)]
  \item Zuordnung:
  \begin{center}
  \begin{tabular}{llll}
  \toprule
  Beschreibung & Begriff & Formel & Skript-Gleichung \\
  \midrule
  A & Konditionierung (conditioning) & (iii) & Gl.~(13) \\
  B & Stabilität (stability) & (iv) & Gl.~(14) \\
  C & Konsistenz (consistency) & (i) & Gl.~(15) \\
  D & Konvergenz (convergence) & (ii) & Gl.~(16) \\
  \bottomrule
  \end{tabular}
  \end{center}
  \item Die \textbf{Konditionierung} ist eine Eigenschaft des \emph{Problems}: In Gl.~(13) kommt nur $f$ vor, nicht $\hat{f}$ --- sie ändert sich nicht, wenn man die Methode wechselt. \textbf{Stabilität, Konsistenz und Konvergenz} sind Eigenschaften der \emph{Methode} (bzw.\ der numerischen Lösung): In den Gln.~(14)--(16) tritt jeweils $\hat{f}$ auf.
  \item In der Konvergenz (Gl.~16) wird $\hat{f}(\tilde{x})$ mit $f(x)$ verglichen. Darin stecken beide Vergleiche zugleich: der Umgang der Methode mit der \emph{gestörten} Eingabe $\tilde{x}$ (das ist die Stabilität, Gl.~14: $\hat{f}(\tilde{x})$ vs.\ $\hat{f}(x)$) und die Nähe der Methode zur \emph{exakten} Lösung (das ist die Konsistenz, Gl.~15: $\hat{f}(x)$ vs.\ $f(x)$). Mit der Dreiecksungleichung gilt entsprechend
  \[ \left|\hat{f}(\tilde{x}) - f(x)\right| \leq \underbrace{\left|\hat{f}(\tilde{x}) - \hat{f}(x)\right|}_{\text{Stabilität}} + \underbrace{\left|\hat{f}(x) - f(x)\right|}_{\text{Konsistenz}} \leq s\,|\tilde{x} - x| + c. \]
  Nur wer Störungen kontrolliert \emph{und} die exakte Lösung approximiert, konvergiert.
\end{enumerate}

\subsection*{Lösung zu Aufgabe~\ref{auf:3}}

\begin{enumerate}[label=(\alph*)]
  \item Ja. Die Fehlerfolge stabilisiert sich für $h \to 0$ sichtbar: $0.080 \to 0.045 \to 0.034 \to 0.0315 \to 0.0311 \to 0.0310$. Die Methode konvergiert im Sinne von Gl.~(16) gegen eine Fehlerschranke von ungefähr $0.031$.
  \item Die Fehlerschranke ist \emph{nicht} null, die Methode konvergiert also gegen einen von $f(x)$ verschiedenen Wert. Nach dem Skript (Randnotiz S.~29) zeigt das einen \textbf{systematischen Fehler (Bias)} der Methode an: \glqq If the method converges to a value different from $f(x)$, this indicates a systematic error (bias) in the method.\grqq{} Kein noch so kleines $h$ beseitigt diesen Restfehler.
  \item \textbf{Bias:} der konstante Restfehler ($\approx 0.031$), der trotz beliebiger Verfeinerung bestehen bleibt --- ein systematischer Versatz der numerischen Lösung gegenüber der Wahrheit (hier z.\,B.\ durch die Eingabestörung $\tilde{x}$ verursacht). \textbf{Varianz:} die Fluktuationen der berechneten Werte, die durch wechselnde Störungen der Eingabedaten entstehen (Streuung der Lösung um ihren Grenzwert, vgl.\ die Instabilität feiner Gitter). Verfeinern von $h$ reduziert den Konsistenzanteil des Fehlers, nicht aber den Bias.
\end{enumerate}

\subsection*{Lösung zu Aufgabe~\ref{auf:4}}

Alle Werte aus der Tabelle; in der letzten Nachkommastelle können Rundungsabweichungen von $\pm 1$ auftreten.

\begin{enumerate}[label=(\alph*)]
  \item Am Punkt $x = 3$ (nahe der Minimalstelle $\pi$, Kosinus flach):
  \begin{align*}
  \kappa(3, +0.25) &= |\cos(3) - \cos(3.25)| = |-0.98999 - (-0.99413)| = 0.00414 \\
  \kappa(3, -0.25) &= |\cos(3) - \cos(2.75)| = |-0.98999 - (-0.92430)| = 0.06569
  \end{align*}
  \item Am Punkt $x = 1.5$ (nahe der Nullstelle $\pi/2$, Kosinus steil):
  \begin{align*}
  \kappa(1.5, +0.25) &= |\cos(1.5) - \cos(1.75)| = |0.07074 - (-0.17825)| = 0.24899 \approx 0.2490 \\
  \kappa(1.5, -0.25) &= |\cos(1.5) - \cos(1.25)| = |0.07074 - 0.31532| = 0.24458 \approx 0.2446
  \end{align*}
  \item Bei $x = 3$ liegt \textbf{gute Konditionierung} (well-conditioning) vor: $0.00414 \leq \kappa \leq 0.06569$ --- die Störung $\pm 0.25$ schlägt nur schwach auf den Funktionswert durch. Bei $x = 1.5$ liegt \textbf{schlechte Konditionierung} (ill-conditioning) vor mit $\kappa \approx 0.249$, also rund dem Vierfachen des ungünstigsten Falls bei $x = 3$ (und zusätzlich ist dort wegen $\cos(1.5) \approx 0$ der \emph{relative} Fehler dramatisch). Begründung über die Steigung: Die lokale Empfindlichkeit wird durch $|f'(x)| = |{-\sin(x)}| = |\sin(x)|$ bestimmt. Bei $x = 3$ ist $|\sin(3)| \approx 0.141$ (flach, nahe der Minimalstelle $\pi$, wo $f' = 0$); bei $x = 1.5$ ist $|\sin(1.5)| \approx 0.997 \approx 1$ (maximale Steigung nahe $\pi/2$). Kleine Eingabeänderungen werden dort also nahezu eins zu eins in Ausgabeänderungen übersetzt.
\end{enumerate}

\subsection*{Lösung zu Aufgabe~\ref{auf:5}}

Betrachtet wird $x = 2$, $\tilde{x} = 2.25$, also $|\tilde{x} - x| = 0.25$.

\begin{enumerate}[label=(\alph*)]
  \item \textbf{Für $h = 1$} (Gitter $\{0, 1, 2, 3, 4\}$): Der zu $x = 2$ nächstgelegene Gitterpunkt ist $2$; der zu $\tilde{x} = 2.25$ nächstgelegene Gitterpunkt ist ebenfalls $2$ (Abstand $0.25$ gegenüber $0.75$ zu $3$). Beide Auswertungen fallen auf denselben Gitterwert:
  \[ \left| \hat{f}_1(2.25) - \hat{f}_1(2) \right| = \left| \cos(2) - \cos(2) \right| = 0 \]
  \[ s_{h=1,\,x=2} = 0 \div 0.25 \approx 0. \]
  Die grobe Diskretisierung ist an diesem Punkt maximal stabil: Sie \glqq verschluckt\grqq{} die Störung vollständig.
  \item \textbf{Für $h = 0.2$} (Gitter $\{0, 0.2, \ldots, 4\}$): Der zu $x = 2$ nächstgelegene Gitterpunkt ist $2.0$; der zu $\tilde{x} = 2.25$ nächstgelegene Gitterpunkt ist $2.2$ (Abstand $0.05$ gegenüber $0.15$ zu $2.4$). Damit:
  \[ \left| \hat{f}_{0.2}(2.25) - \hat{f}_{0.2}(2) \right| = \left| \cos(2.2) - \cos(2) \right| = |-0.58850 - (-0.41615)| = 0.17235 \]
  \[ s_{h=0.2,\,x=2} = 0.17235 \div 0.25 \approx 0.6894. \]
  \item Die \textbf{grobe Diskretisierung ($h = 1$) ist stabiler} ($s \approx 0$ gegenüber $s \approx 0.6894$). Das feine Gitter \glqq sieht\grqq{} die Störung: $x$ und $\tilde{x}$ fallen auf \emph{verschiedene} Gitterpunkte mit verschiedenen Funktionswerten, sodass die Methode auf die Eingabestörung reagiert. Beim groben Gitter landen $x$ und $\tilde{x}$ auf demselben Gitterpunkt --- die Ausgabe ändert sich gar nicht. (Vgl.\ die Warnung des Skripts, S.~31: In Randregionen der Diskretisierung hat die Stabilität Anomalien, besonders bei stückweise konstanten Approximationen --- läge $\tilde{x}$ jenseits der Zellgrenze, spränge auch das grobe Gitter.)
\end{enumerate}

\subsection*{Lösung zu Aufgabe~\ref{auf:6}}

\begin{enumerate}[label=(\alph*)]
  \item \textbf{Konsistenz} an $x^* = \pi \approx 3.14159$, $f(x^*) = -1$.

  Für $h = 1$: Der nächstgelegene Gitterpunkt zu $\pi$ ist $3$ (Abstand $0.14159$ gegenüber $0.85841$ zu $4$):
  \[ c_1 \geq \left| \hat{f}_1(\pi) - f(\pi) \right| = \left| \cos(3) - (-1) \right| = |-0.98999 + 1| = 0.01001. \]
  Für $h = 0.2$: Der nächstgelegene Gitterpunkt zu $\pi$ ist $3.2$ (Abstand $0.05841$ gegenüber $0.14159$ zu $3.0$):
  \[ c_{0.2} \geq \left| \hat{f}_{0.2}(\pi) - f(\pi) \right| = \left| \cos(3.2) - (-1) \right| = |-0.99829 + 1| = 0.00171. \]
  Die Konsistenz verbessert sich erwartungsgemäß mit kleinerer Schrittweite: $c_1 = 0.01001$ gegenüber $c_{0.2} = 0.00171$.

  \item \textbf{Konvergenz} mit $\tilde{x} = \pi + 0.25 \approx 3.39159$.

  Für $h = 1$: Der nächstgelegene Gitterpunkt zu $3.39159$ ist $3$ (Abstand $0.39159$ gegenüber $0.60841$ zu $4$):
  \[ \left| \hat{f}_1(\tilde{x}) - f(x^*) \right| = \left| \cos(3) - (-1) \right| = |-0.98999 + 1| = 0.01001. \]
  Für $h = 0.2$: Der nächstgelegene Gitterpunkt zu $3.39159$ ist $3.4$ (Abstand $0.00841$ gegenüber $0.19159$ zu $3.2$):
  \[ \left| \hat{f}_{0.2}(\tilde{x}) - f(x^*) \right| = \left| \cos(3.4) - (-1) \right| = |-0.96680 + 1| = 0.03320. \]

  \item \textbf{Befund:} Beim Übergang $h = 1 \to h = 0.2$ verbessert sich die Konsistenz ($0.01001 \to 0.00171$), aber die Konvergenz \emph{verschlechtert} sich ($0.01001 \to 0.03320$): Das feine Gitter folgt der gestörten Eingabe (Instabilität, vgl.\ Aufgabe~5) und trägt deren Fehler voll in die Ausgabe, während das grobe Gitter die Störung wegmittelt. Eine feinere Diskretisierung ist also \emph{nicht} automatisch besser. (Zum Vergleich: Im Sinus-Beispiel des Skripts lieferten beide Schrittweiten denselben Konvergenzfehler $0.15853$ [Korrektur ggü.\ Original, S.~32: das PDF druckt dort \glqq$0.1599$\grqq; korrekt ist $|-0.84147 - (-1)| = 0.15853$] --- auch dort verbesserte kleineres $h$ die Konvergenz nicht.)
\end{enumerate}

\subsection*{Lösung zu Aufgabe~\ref{auf:7}}

\begin{enumerate}[label=(\alph*)]
  \item Sei $x \in [a,b]$ beliebig und $x_i$ der zu $x$ nächstgelegene Gitterpunkt. Da die Gitterpunkte den Abstand $h$ haben, gilt für den nächstgelegenen Gitterpunkt stets
  \[ |x - x_i| \leq \frac{h}{2}. \]
  Nach dem Mittelwertsatz (bzw.\ der Lipschitz-Stetigkeit von $f$ mit Konstante $L = \max |f'|$) existiert ein $\xi$ zwischen $x$ und $x_i$ mit
  \[ \left| \hat{f}_h(x) - f(x) \right| = \left| f(x_i) - f(x) \right| = |f'(\xi)| \cdot |x_i - x| \leq L \cdot \frac{h}{2}. \]
  Für $h \to 0$ geht die rechte Seite gegen null --- die Methode wird punktweise exakt, gleichmäßig in $x$.
  \item Mit der Dreiecksungleichung (zweimal angewendet, \glqq umgekehrte\grqq{} Dreiecksungleichung) folgt
  \[ \Big| \left| \hat{f}_h(\tilde{x}) - \hat{f}_h(x) \right| - \left| f(\tilde{x}) - f(x) \right| \Big|
     \leq \left| \hat{f}_h(\tilde{x}) - f(\tilde{x}) \right| + \left| \hat{f}_h(x) - f(x) \right|
     \leq 2 \cdot L \cdot \frac{h}{2} = L \, h \xrightarrow{\;h \to 0\;} 0, \]
  also
  \[ \left| \hat{f}_h(\tilde{x}) - \hat{f}_h(x) \right| \xrightarrow{\;h \to 0\;} \left| f(\tilde{x}) - f(x) \right| = \kappa. \]
  \textbf{Interpretation:} Die linke Seite der Stabilitätsungleichung (Gl.~14) --- die Empfindlichkeit der \emph{Methode} gegenüber der Eingabestörung --- geht im Grenzfall $h \to 0$ exakt in die Konditionszahl $\kappa$ --- die Empfindlichkeit des \emph{Problems} --- über. Die kleinstmögliche Stabilitätskonstante konvergiert entsprechend gegen
  \[ s_h^{\min} = \frac{\left| \hat{f}_h(\tilde{x}) - \hat{f}_h(x) \right|}{|\tilde{x} - x|} \xrightarrow{\;h \to 0\;} \frac{\kappa}{|\tilde{x} - x|}. \]
  Für $h \to 0$ erbt die Methode also die Empfindlichkeit des Problems: \emph{Stabilität = Konditionierung}. Stabiler als das Problem es zulässt, kann eine (punktweise exakte) Methode nicht werden.
  \item Numerische Überprüfung bei $x = 2$, $\tilde{x} = 2.25$:
  \begin{center}
  \begin{tabular}{lccc}
  \toprule
   & $h = 1$ & $h = 0.2$ & $h = 0.05$ \\
  \midrule
  $\left| \hat{f}_h(2.25) - \hat{f}_h(2) \right|$ & $0$ & $0.17235$ & $0.21202$ \\
  \bottomrule
  \end{tabular}
  \end{center}
  Für $h = 0.05$ sind $x = 2$ und $\tilde{x} = 2.25$ selbst Gitterpunkte, daher $\left|\hat{f}_{0.05}(2.25) - \hat{f}_{0.05}(2)\right| = |\cos(2.25) - \cos(2)| = |-0.62817 + 0.41615| = 0.21202$. Die Konditionszahl ist
  \[ \kappa(2, +0.25) = |\cos(2) - \cos(2.25)| = 0.21202. \]
  Die Folge $0 \to 0.17235 \to 0.21202$ nähert sich mit fallendem $h$ der Konditionszahl an und erreicht sie hier exakt --- wie behauptet.
\end{enumerate}

\subsection*{Lösung zu Aufgabe~\ref{auf:8}}

\begin{enumerate}[label=(\alph*)]
  \item Beim Übergang $h = 1 \to h = 0.2$ wird die \textbf{Konsistenz} besser ($c$: $0.01001 \to 0.00171$ --- der reine Diskretisierungsfehler sinkt mit $h$), die \textbf{Stabilität} aber schlechter ($s$: $\approx 0 \to \approx 0.6894$ --- das feine Gitter reagiert auf die Eingabestörung). Da die \textbf{Konvergenz} beide Anteile kombiniert ($|\hat{f}(\tilde{x}) - f(x)| \leq s\,|\tilde{x}-x| + c$, vgl.\ Aufgabe~2\,c), kann der Gesamtfehler trotz besserer Konsistenz wachsen --- hier von $0.01001$ auf $0.03320$. Bessere Konsistenz allein garantiert also keine bessere Konvergenz, weil der über die Stabilität eingeschleppte Störungsanteil dominieren kann.
  \item Nach Aufgabe~7 wird $\hat{f}_h$ für $h \to 0$ punktweise exakt, also
  \[ \left| \hat{f}_h(\tilde{x}) - f(x^*) \right| \xrightarrow{\;h \to 0\;} \left| f(\tilde{x}) - f(x^*) \right| = \left| \cos(\pi + 0.25) - (-1) \right| = |-0.96891 + 1| = 0.03109. \]
  Der Konvergenzfehler konvergiert also gegen $\approx 0.0311 \neq 0$: \textbf{nicht verschwindende Konvergenz (non-zero convergence)}. Wie in Aufgabe~3 zeigt der von null verschiedene Grenzwert einen \textbf{systematischen Fehler (Bias)} an --- hier verursacht durch die Eingabestörung $\tilde{x} = x^* + 0.25$, die keine noch so feine Diskretisierung beheben kann. (Genau dieser Wert ist die Konditionsschranke des gestörten Problems an der Minimalstelle.)
  \item Zwei unabhängige Gründe:
  \begin{itemize}[nosep]
    \item \textbf{Konditionierung der gestörten Eingabe (Bias-Schranke):} Solange die Eingabe um $\pm 0.25$ gestört ist, bleibt nach (b) ein Restfehler der Größenordnung $\kappa$ bestehen --- die Methode kann das Problem nicht besser lösen, als seine Konditionierung es zulässt. Verkleinern von $h$ senkt nur den Konsistenzanteil, nicht den Störungsanteil.
    \item \textbf{Gleitkommadarstellung (floating-point representation):} Zahlen werden mit endlicher Präzision gespeichert; jede Operation kann einen Rundungsfehler der relativen Größe $\varepsilon_{\text{mach}}$ einschleppen. Kleineres $h$ bedeutet mehr Gitterpunkte ($N = (b-a)/h$) und damit mehr Operationen, deren Rundungsfehler akkumulieren und sich verstärken können. Zudem unterscheiden sich benachbarte Funktionswerte nur noch um etwa $|f'| \cdot h$; fällt dieser Abstand unter die absolute Auflösung $\varepsilon_{\text{mach}} \cdot |f(x_i)|$, sind Nachbarwerte numerisch nicht mehr unterscheidbar, und bei Differenzen nahezu gleicher Werte droht Auslöschung (loss of significance / catastrophic cancellation). Der Gesamtfehler verläuft daher typisch U-förmig über $h$: Der Abschneide-/Konsistenzfehler fällt, der Rundungsfehler steigt --- unterhalb eines optimalen $h$ wird das Ergebnis wieder schlechter. (Hinzu kommt als praktisches Problem der wachsende Rechenaufwand.)
  \end{itemize}
\end{enumerate}

\end{document}
